%% ma: L12-B05-0016
%% lop: 12
%% chuong: 1
%% bai: 5
%% dang: 12.5.11
%% loai: TLN
%% muc: VD
%% dap_an: "60"
%% nguon: "Ảnh giáo viên gửi 11/10/2026 (Đề kiểm tra giữa HK1 Toán 12, THPT Việt Đức 2026-2027, mã đề 0103), Phần III Câu 2"
%% kien_thuc: "Bài 5 (bài toán tối ưu bằng đạo hàm)"
%% trang_thai: cho-duyet
%% ly_do: "Dạng mới đề xuất, chờ giáo viên duyệt."
\begin{ex}
Một xe điện giao hàng đi quãng đường bằng phẳng dài $120$ km với vận tốc không đổi $v$ (km/h). Công suất tiêu thụ điện khi xe chạy ở vận tốc này được mô hình hóa bởi hàm số $P(v)=12+\dfrac{v^2}{300}$ (kW). Bỏ qua điện năng lúc khởi hành và dừng xe. Để điện năng tiêu thụ trên cả quãng đường là nhỏ nhất, xe cần chạy với vận tốc bao nhiêu kilômét trên giờ? (Biết công thức tính điện năng tiêu thụ (kWh) là: điện năng $=$ công suất tiêu thụ điện $\times$ thời gian.)
\shortans{60}
\loigiai{Thời gian đi hết quãng đường là $t=\dfrac{120}v$ (giờ) nên điện năng tiêu thụ là
\[E(v)=\left(12+\dfrac{v^2}{300}\right)\dfrac{120}v=\dfrac{1440}v+\dfrac{2v}5,\quad v>0.\]
$E'(v)=-\dfrac{1440}{v^2}+\dfrac25=0\Leftrightarrow v^2=3600\Leftrightarrow v=60$. $E'(v)<0$ trên $(0;60)$ và $E'(v)>0$ trên $(60;+\infty)$ nên $E$ nhỏ nhất tại $v=60$ ($E(60)=48$ kWh). Vậy xe cần chạy với vận tốc $60$ km/h.}
\end{ex}
